% ------------------------------------------------------------------
% Algebra Basic Skills booklet.
% GENERATED by src/build.py from src/content.py -- edit content.py and
% re-run the build, so every practice answer gets re-checked.
% ------------------------------------------------------------------
\documentclass[12pt,a4paper]{article}
\usepackage[a4paper,top=1.4cm,bottom=1.8cm,left=1.5cm,right=1.5cm,footskip=0.8cm]{geometry}
\usepackage[T1]{fontenc}
\usepackage[scaled=0.95]{helvet}
\renewcommand{\familydefault}{\sfdefault}
\usepackage{amsmath,amssymb}
\usepackage[table]{xcolor}
\usepackage[most]{tcolorbox}
\usepackage{enumitem,multicol,array,tabularx,pifont,fancyhdr,titlesec}

\definecolor{accent}{HTML}{1D6B8C}
\definecolor{accentlight}{HTML}{E3F0F6}
\definecolor{checkcol}{HTML}{8A3B12}
\definecolor{checklight}{HTML}{F6E7DE}
\definecolor{muted}{HTML}{56616B}
\definecolor{good}{HTML}{2E7D32}
\definecolor{warm}{HTML}{FFF6DD}
\definecolor{warmedge}{HTML}{E3B341}

\setlength{\parindent}{0pt}
\setlength{\parskip}{4pt}
\newcommand{\cmark}{\textcolor{good}{\ding{51}}}
\newcommand{\xmark}{\textcolor{red!70!black}{\ding{55}}}
\newcommand{\tickbox}{{\Large$\square$}}
\titleformat*{\subsection}{\large\bfseries\color{accent}}
\titlespacing*{\subsection}{0pt}{8pt}{4pt}

\pagestyle{fancy}
\fancyhf{}
\renewcommand{\headrulewidth}{0pt}
\fancyfoot[L]{\footnotesize\color{muted}Algebra Basic Skills}
\fancyfoot[R]{\footnotesize\color{muted}Page \thepage}

% --- boxes -----------------------------------------------------------------
\tcbset{base/.style={enhanced, arc=2mm, boxrule=0.8pt, left=3mm, right=3mm, top=1.5mm, bottom=1.5mm}}
\newtcolorbox{tipbox}[1]{base, fontupper=\large, colframe=accent, colback=white, before upper={\textbf{#1.}\ }}
\newtcolorbox{stepsbox}[1]{base, fontupper=\large, colframe=accent, colback=white, before upper={\textbf{#1}\par}}
\newtcolorbox{rememberbox}{base, fontupper=\large, colframe=warmedge, colback=warm, before upper={\textbf{Remember:}\ }}
\newenvironment{examples}{\begin{tcbraster}[raster columns=3, raster equal height=rows,
  raster column skip=2.5mm, raster row skip=2.5mm, raster before skip=2mm, raster after skip=3mm]}{\end{tcbraster}}
\newcommand{\example}[3]{\begin{tcolorbox}[base, colframe=black!25, colback=white, boxrule=0.5pt, arc=1.5mm]
  {\Large #1}\par\vspace{-2pt}{\color{black!25}\rule{\linewidth}{0.4pt}}\par
  {\normalsize\color{black!75}#2}\par\vspace{2pt}
  {\Large\color{good}\boldmath #3}\end{tcolorbox}}

\newcommand{\stepnum}[2]{\tikz[baseline=(n.base)]\node[circle,fill=#2,text=white,inner sep=0pt,
  minimum size=8mm,font=\large\bfseries](n){#1};}
\newcommand{\qnum}[1]{\tikz[baseline=(n.base)]\node[circle,fill=accentlight,text=accent,inner sep=0pt,
  minimum size=6mm,font=\small\bfseries](n){#1};}
\newcommand{\cnum}[1]{\tikz[baseline=(n.base)]\node[circle,fill=checklight,text=checkcol,inner sep=0pt,
  minimum size=6.5mm,font=\small\bfseries](n){#1};}

% \banner{number}{title}{label}{colour}
\newcommand{\banner}[4]{\begin{tcolorbox}[enhanced, colback=#4, colframe=#4, arc=2mm,
  left=3mm, right=3mm, top=2mm, bottom=2mm, after skip=3mm]
  \color{white}\tikz[baseline=(n.base)]\node[circle,fill=white,text=#4,inner sep=0pt,
  minimum size=8.5mm,font=\Large\bfseries](n){#1};\hspace{3mm}{\LARGE\bfseries #2}\hfill
  \colorbox{white!20!#4}{\bfseries\MakeUppercase{#3}}\end{tcolorbox}}

% \qbox{n}{question}{height}
\newcommand{\qbox}[3]{\begin{tcolorbox}[base, height=#3, text fill, colframe=black!30, colback=white, boxrule=0.5pt,
  arc=1.5mm, top=1.5mm, bottom=2mm]\qnum{#1}\hspace{2.5mm}{\Large #2}\par\vfill
  {\Large$=$}\hspace{2mm}\textcolor{black!45}{\rule{0.8\linewidth}{0.6pt}}\end{tcolorbox}}

% \cbox{n}{question}{height}  -- skill check question, with a mark box
\newcommand{\cbox}[3]{\begin{tcolorbox}[base, height=#3, text fill, colframe=checkcol!70, colback=white, boxrule=0.8pt,
  arc=1.5mm, top=2mm, bottom=2.5mm]\cnum{#1}\hspace{2.5mm}{\Large #2}\hfill
  \fcolorbox{checkcol}{white}{\rule{0pt}{6mm}\hspace{6mm}}\par\vfill
  {\Large$=$}\hspace{2mm}\textcolor{black!45}{\rule{0.8\linewidth}{0.6pt}}\end{tcolorbox}}

\newcommand{\fillline}[1]{\textcolor{black}{\rule{#1}{0.6pt}}}

\begin{document}

\thispagestyle{empty}
\begin{tcolorbox}[enhanced, colback=accent, colframe=accent, arc=3mm, halign=center, top=9mm, bottom=8mm]
\color{white}{\small\bfseries B\,A\,S\,I\,C\quad S\,K\,I\,L\,L\,S\quad B\,O\,O\,K\,L\,E\,T}\par\vspace{3mm}
{\fontsize{34}{40}\selectfont\bfseries Algebra Basic Skills}\par\vspace{3mm}
{\Large Learn it. Practise it. Show you can do it.}
\end{tcolorbox}
\vspace{6mm}
{\large Name: \fillline{9.5cm} \hfill Class: \fillline{3.5cm}}

\vspace{4mm}
{\Large\bfseries\color{accent} How to use this booklet}
\begin{tcbraster}[raster columns=1, raster row skip=2mm]
\begin{tcolorbox}[base, colback=accentlight, colframe=accentlight]
\stepnum{1}{accent}\hspace{3mm}
\begin{minipage}[t]{0.88\linewidth}\textbf{Learn.} Read the notes and the worked examples.
Each example shows every step.\end{minipage}
\end{tcolorbox}
\begin{tcolorbox}[base, colback=accentlight, colframe=accentlight]
\stepnum{2}{accent}\hspace{3mm}
\begin{minipage}[t]{0.88\linewidth}\textbf{Practise.} Do \textbf{Practice A}, then \textbf{Practice B}.
Check your answers at the back of the booklet and fix any mistakes.\end{minipage}
\end{tcolorbox}
\begin{tcolorbox}[base, colback=checklight, colframe=checklight]
\stepnum{3}{checkcol}\hspace{3mm}
\begin{minipage}[t]{0.88\linewidth}\textbf{Skill Check.} When you're ready, do the Skill Check
\textbf{on your own, without notes}. Your teacher will mark it.\end{minipage}
\end{tcolorbox}
\end{tcbraster}

\vspace{3mm}
{\Large\bfseries\color{accent} My progress}\par\vspace{1mm}
\renewcommand{\arraystretch}{1.9}
\begin{tabularx}{\linewidth}{|c|X|c|c|p{2.2cm}|}
\hline
\rowcolor{accentlight} & \textbf{Skill} & \small\textbf{Practice A} & \small\textbf{Practice B} & \small\textbf{Skill Check} \\
\rowcolor{accentlight} & & \small done & \small done & \small (teacher) \\ \hline
\textbf{\color{accent}\large 1} & \textbf{Collecting Like Terms}\newline {\small\color{muted}I can add and subtract like terms, e.g. $3a + 4a = 7a$} & \tickbox & \tickbox & \hfill/\,4 \\ \hline
\textbf{\color{accent}\large 2} & \textbf{Multiplying Algebra}\newline {\small\color{muted}I can multiply terms, e.g. $2a \times  3b = 6ab$ and $x \times  x = x^{2}$} & \tickbox & \tickbox & \hfill/\,4 \\ \hline
\textbf{\color{accent}\large 3} & \textbf{Dividing Algebra}\newline {\small\color{muted}I can divide terms, e.g. $8x \div  2 = 4x$ and $6ab \div  2a = 3b$} & \tickbox & \tickbox & \hfill/\,4 \\ \hline
\textbf{\color{accent}\large 4} & \textbf{Expanding Brackets}\newline {\small\color{muted}I can expand brackets, e.g. $3(x + 2) = 3x + 6$ and $-2(x - 4) = -2x + 8$} & \tickbox & \tickbox & \hfill/\,4 \\ \hline
\textbf{\color{accent}\large 5} & \textbf{Factorising into Brackets}\newline {\small\color{muted}I can factorise, e.g. $2x + 6 = 2(x + 3)$} & \tickbox & \tickbox & \hfill/\,4 \\ \hline
& \textbf{Mixed Review} \small(extra practice) & \multicolumn{2}{c|}{\tickbox} & \textbf{Total}\hfill/\,20 \\ \hline
\end{tabularx}

\vfill
{\small\color{muted} The answers to the \textbf{practice} questions are at the back.
The Skill Checks have no answers in this booklet --- that's where you show what you can do!}
\newpage
\banner{1}{Collecting Like Terms}{Learn}{accent}
\begin{tcolorbox}[enhanced, colback=accentlight, colframe=accent, boxrule=0pt, leftrule=1.5mm, arc=0mm, top=1.5mm, bottom=1.5mm]\textbf{Goal:} I can add and subtract like terms, e.g. $3a + 4a = 7a$\end{tcolorbox}
\large

\textbf{Like terms} have \textbf{exactly the same letter part}. You can only add or
subtract like terms.

\medskip
\begin{center}
\renewcommand{\arraystretch}{1.35}
\begin{tabular}{|l|l|}
\hline
\rowcolor{accentlight}\textbf{Like terms \cmark} & \textbf{NOT like terms \xmark} \\ \hline
$3a$ and $5a$ & $3a$ and $5b$ \quad (different letters) \\ \hline
$x$ and $7x$ & $x$ and $7$ \quad (one has no letter) \\ \hline
$2x^{2}$ and $x^{2}$ & $x^{2}$ and $x$ \quad ($x^{2}$ is not the same as $x$) \\ \hline
$4ab$ and $ab$ & $ab$ and $a$ \\ \hline
\end{tabular}
\end{center}

\begin{tipbox}{Think of the letters as objects}
3 apples + 4 apples = 7 apples, so $3a + 4a = 7a$.\\
But 3 apples + 2 bananas can't be joined, so $3a + 2b$ stays as $3a + 2b$.
\end{tipbox}

\begin{tipbox}{A letter on its own means 1 of it}
$a$ is the same as $1a$. So $a + a + a = 3a$.
\end{tipbox}

\subsection*{Worked examples}
\begin{examples}
\example{$3a + 4a$}{3 lots of $a$ plus 4 lots of $a$}{$= 7a$}
\example{$6x - 2x$}{6 take away 2 is 4}{$= 4x$}
\example{$5y - y$}{$y$ means $1y$, and $5 - 1 = 4$}{$= 4y$}
\example{$2a + 3b + 4a + b$}{The $a$'s: $2a + 4a = 6a$\\The $b$'s: $3b + b = 4b$}{$= 6a + 4b$}
\example{$x^{2} + 3x + 2x^{2}$}{The $x^{2}$'s: $x^{2} + 2x^{2} = 3x^{2}$\\The $x$'s: just $3x$}{$= 3x^{2} + 3x$}
\example{$7m + 2 - 3m$}{The $m$'s: $7m - 3m = 4m$\\The numbers: just $2$}{$= 4m + 2$}
\end{examples}

\begin{rememberbox}
The sign ($+$ or $-$) in \textbf{front} of a term belongs to that term.
Circle or highlight each type of term in a different colour before you add them up.
\end{rememberbox}

\newpage\normalsize
\banner{1}{Collecting Like Terms}{Practice A}{accent}
{\large\textbf{Simplify by collecting like terms.}} {\color{muted}Show your working in each box.}\par\vspace{1mm}
\begin{tcbraster}[raster columns=2, raster column skip=4mm, raster row skip=3mm]
\qbox{1}{$a + a$}{3.45cm}
\qbox{2}{$a + a + a$}{3.45cm}
\qbox{3}{$2a + 3a$}{3.45cm}
\qbox{4}{$4x + 5x$}{3.45cm}
\qbox{5}{$3m + m$}{3.45cm}
\qbox{6}{$7b - 3b$}{3.45cm}
\qbox{7}{$6y - y$}{3.45cm}
\qbox{8}{$10x - 4x$}{3.45cm}
\qbox{9}{$2p + 3p + p$}{3.45cm}
\qbox{10}{$8c - 2c$}{3.45cm}
\qbox{11}{$x^{2} + x^{2}$}{3.45cm}
\qbox{12}{$3x^{2} + 4x^{2}$}{3.45cm}
\end{tcbraster}
\vfill\centerline{\color{good}\bfseries \ding{51}\ Check your answers at the back. Fix any you got wrong before moving on.}
\newpage
\banner{1}{Collecting Like Terms}{Practice B}{accent}
{\large\textbf{Simplify by collecting like terms. Some have two types of term.}} {\color{muted}Show your working in each box.}\par\vspace{1mm}
\begin{tcbraster}[raster columns=2, raster column skip=4mm, raster row skip=3mm]
\qbox{1}{$2a + 3b + 4a$}{3.45cm}
\qbox{2}{$5x + 2y + x$}{3.45cm}
\qbox{3}{$3a + 2b + a + 4b$}{3.45cm}
\qbox{4}{$4x + 3 + 2x$}{3.45cm}
\qbox{5}{$7m - 2m + 3n$}{3.45cm}
\qbox{6}{$5a + 2b - 3a$}{3.45cm}
\qbox{7}{$6x + 4 - 2x + 1$}{3.45cm}
\qbox{8}{$8x - 3x + 2y$}{3.45cm}
\qbox{9}{$4ab + 2ab$}{3.45cm}
\qbox{10}{$x^{2} + 2x + 3x$}{3.45cm}
\qbox{11}{$2x^{2} + x + x^{2}$}{3.45cm}
\qbox{12}{$3y^{2} + 2y + y^{2} + 4y$}{3.45cm}
\end{tcbraster}
\vfill\centerline{\color{good}\bfseries \ding{51}\ Check your answers at the back. Fix any you got wrong before moving on.}
\newpage
\banner{1}{Collecting Like Terms}{Skill Check}{checkcol}
{\large Name: \fillline{8cm}}\hfill
\fcolorbox{checkcol}{white}{\parbox[c][9mm][c]{4.6cm}{\centering\color{checkcol}\bfseries\large Teacher mark: \underline{\hspace{9mm}} / 4}}
\par\vspace{2mm}
{\large\textbf{Simplify by collecting like terms.}} Do this on your own --- no notes. Show your working, and write your final answer on the line.
\par\vspace{2mm}
\begin{tcbraster}[raster columns=2, raster column skip=5mm, raster row skip=5mm]
\cbox{1}{$5a + 2a$}{10.6cm}
\cbox{2}{$8x - 3x$}{10.6cm}
\cbox{3}{$4a + 3b + 2a$}{10.6cm}
\cbox{4}{$x^{2} + 5x + 2x^{2}$}{10.6cm}
\end{tcbraster}
\newpage
\banner{2}{Multiplying Algebra}{Learn}{accent}
\begin{tcolorbox}[enhanced, colback=accentlight, colframe=accent, boxrule=0pt, leftrule=1.5mm, arc=0mm, top=1.5mm, bottom=1.5mm]\textbf{Goal:} I can multiply terms, e.g. $2a \times  3b = 6ab$ and $x \times  x = x^{2}$\end{tcolorbox}
\large

In algebra we \textbf{don't write the $\times$ sign}. Letters and numbers that are multiplied
are written next to each other, \textbf{number first}.

\medskip
\begin{center}
\renewcommand{\arraystretch}{1.35}
\begin{tabular}{|c|l|}
\hline
\rowcolor{accentlight}\textbf{We write\ldots} & \textbf{\ldots it means} \\ \hline
$3a$ & $3 \times  a$ \\ \hline
$ab$ & $a \times  b$ \\ \hline
$5xy$ & $5 \times  x \times  y$ \\ \hline
$x^{2}$ & $x \times  x$ \quad (say \textquotedblleft{}$x$ squared'') \\ \hline
\end{tabular}
\end{center}

\begin{stepsbox}{How to multiply terms}
\begin{enumerate}[nosep]
\item \textbf{Multiply the numbers} together.
\item \textbf{Write the letters} after the number, in alphabetical order.
\item If the \textbf{same letter} appears twice, write it as a \textbf{square}: $x \times  x = x^{2}$.
\end{enumerate}
\end{stepsbox}

\subsection*{Worked examples}
\begin{examples}
\example{$4 \times  a$}{Number first, then the letter}{$= 4a$}
\example{$a \times  b$}{Write the letters together}{$= ab$}
\example{$2a \times  4$}{Numbers: $2 \times  4 = 8$\\Letters: $a$}{$= 8a$}
\example{$2a \times  3b$}{Numbers: $2 \times  3 = 6$\\Letters: $a$ and $b$}{$= 6ab$}
\example{$x \times  x$}{The same letter twice is \textquotedblleft{}squared''}{$= x^{2}$}
\example{$3x \times  2x$}{Numbers: $3 \times  2 = 6$\\Letters: $x \times  x = x^{2}$}{$= 6x^{2}$}
\end{examples}

\begin{rememberbox}
When you \textbf{multiply}, the numbers multiply ($3 \times  2 = 6$) --- don't add them!
And $x \times  x$ is $x^{2}$, not $2x$.\\[2pt]
\textbf{Going further:} $x \times  x \times  x = x^{3}$ (\textquotedblleft{}$x$ cubed''), and $x^{2} \times  x = x^{3}$.
\end{rememberbox}

\newpage\normalsize
\banner{2}{Multiplying Algebra}{Practice A}{accent}
{\large\textbf{Simplify.}} {\color{muted}Show your working in each box.}\par\vspace{1mm}
\begin{tcbraster}[raster columns=2, raster column skip=4mm, raster row skip=3mm]
\qbox{1}{$4 \times  a$}{3.45cm}
\qbox{2}{$5 \times  x$}{3.45cm}
\qbox{3}{$a \times  b$}{3.45cm}
\qbox{4}{$x \times  y$}{3.45cm}
\qbox{5}{$2 \times  3a$}{3.45cm}
\qbox{6}{$3x \times  4$}{3.45cm}
\qbox{7}{$3 \times  5m$}{3.45cm}
\qbox{8}{$6 \times  2x$}{3.45cm}
\qbox{9}{$a \times  a$}{3.45cm}
\qbox{10}{$y \times  y$}{3.45cm}
\qbox{11}{$2a \times  b$}{3.45cm}
\qbox{12}{$p \times  q \times  r$}{3.45cm}
\end{tcbraster}
\vfill\centerline{\color{good}\bfseries \ding{51}\ Check your answers at the back. Fix any you got wrong before moving on.}
\newpage
\banner{2}{Multiplying Algebra}{Practice B}{accent}
{\large\textbf{Simplify. Watch out for the same letter twice!}} {\color{muted}Show your working in each box.}\par\vspace{1mm}
\begin{tcbraster}[raster columns=2, raster column skip=4mm, raster row skip=3mm]
\qbox{1}{$2a \times  3b$}{3.45cm}
\qbox{2}{$4x \times  2y$}{3.45cm}
\qbox{3}{$5a \times  2b$}{3.45cm}
\qbox{4}{$2ab \times  3$}{3.45cm}
\qbox{5}{$3a \times  a$}{3.45cm}
\qbox{6}{$2x \times  5x$}{3.45cm}
\qbox{7}{$4m \times  3m$}{3.45cm}
\qbox{8}{$3x \times  2x$}{3.45cm}
\qbox{9}{$a \times  2b \times  3$}{3.45cm}
\qbox{10}{$x \times  x \times  x$}{3.45cm}
\qbox{11}{$x^{2} \times  x$}{3.45cm}
\qbox{12}{$2x^{2} \times  3x$}{3.45cm}
\end{tcbraster}
\vfill\centerline{\color{good}\bfseries \ding{51}\ Check your answers at the back. Fix any you got wrong before moving on.}
\newpage
\banner{2}{Multiplying Algebra}{Skill Check}{checkcol}
{\large Name: \fillline{8cm}}\hfill
\fcolorbox{checkcol}{white}{\parbox[c][9mm][c]{4.6cm}{\centering\color{checkcol}\bfseries\large Teacher mark: \underline{\hspace{9mm}} / 4}}
\par\vspace{2mm}
{\large\textbf{Simplify.}} Do this on your own --- no notes. Show your working, and write your final answer on the line.
\par\vspace{2mm}
\begin{tcbraster}[raster columns=2, raster column skip=5mm, raster row skip=5mm]
\cbox{1}{$3 \times  5a$}{10.6cm}
\cbox{2}{$2a \times  4b$}{10.6cm}
\cbox{3}{$x \times  x$}{10.6cm}
\cbox{4}{$3x \times  2x$}{10.6cm}
\end{tcbraster}
\newpage
\banner{3}{Dividing Algebra}{Learn}{accent}
\begin{tcolorbox}[enhanced, colback=accentlight, colframe=accent, boxrule=0pt, leftrule=1.5mm, arc=0mm, top=1.5mm, bottom=1.5mm]\textbf{Goal:} I can divide terms, e.g. $8x \div  2 = 4x$ and $6ab \div  2a = 3b$\end{tcolorbox}
\large

Dividing is the opposite of multiplying. A division can also be written as a fraction:
$6a \div  2$ is the same as $\dfrac{6a}{2}$.

\begin{stepsbox}{How to divide terms}
\begin{enumerate}[nosep]
\item \textbf{Divide the numbers.}
\item \textbf{Cancel any letter} that is on the top \textbf{and} the bottom.
\item Write down what is left.
\end{enumerate}
\end{stepsbox}

\begin{tipbox}{Why can letters cancel?}
Anything divided by itself is 1. Just like $5 \div  5 = 1$, $a \div  a = 1$.
So $\dfrac{ab}{a} = b$.
\end{tipbox}

\subsection*{Worked examples}
\begin{examples}
\example{$6a \div  2$}{Numbers: $6 \div  2 = 3$\\Letters: $a$ stays}{$= 3a$}
\example{$ab \div  a$}{The $a$ cancels\\$b$ is left}{$= b$}
\example{$8x \div  x$}{The $x$ cancels\\$8$ is left}{$= 8$}
\example{$12ab \div  4a$}{Numbers: $12 \div  4 = 3$\\$a$ cancels, $b$ is left}{$= 3b$}
\example{$x^{2} \div  x$}{$x^{2}$ is $x \times  x$\\One $x$ cancels, one is left}{$= x$}
\example{$6x^{2} \div  3x$}{Numbers: $6 \div  3 = 2$\\Letters: $x^{2} \div  x = x$}{$= 2x$}
\end{examples}

\begin{rememberbox}
If \textbf{all} the letters cancel, you are left with just a number, e.g. $5a \div  a = 5$.
Write $x^{2}$ as $x \times  x$ if it helps you see what cancels.
\end{rememberbox}

\newpage\normalsize
\banner{3}{Dividing Algebra}{Practice A}{accent}
{\large\textbf{Simplify.}} {\color{muted}Show your working in each box.}\par\vspace{1mm}
\begin{tcbraster}[raster columns=2, raster column skip=4mm, raster row skip=3mm]
\qbox{1}{$8a \div  2$}{3.45cm}
\qbox{2}{$12x \div  3$}{3.45cm}
\qbox{3}{$10b \div  5$}{3.45cm}
\qbox{4}{$20y \div  4$}{3.45cm}
\qbox{5}{$9x \div  3$}{3.45cm}
\qbox{6}{$15p \div  5$}{3.45cm}
\qbox{7}{$18x \div  9$}{3.45cm}
\qbox{8}{$6m \div  6$}{3.45cm}
\qbox{9}{$ab \div  b$}{3.45cm}
\qbox{10}{$xy \div  x$}{3.45cm}
\qbox{11}{$5a \div  a$}{3.45cm}
\qbox{12}{$7x \div  x$}{3.45cm}
\end{tcbraster}
\vfill\centerline{\color{good}\bfseries \ding{51}\ Check your answers at the back. Fix any you got wrong before moving on.}
\newpage
\banner{3}{Dividing Algebra}{Practice B}{accent}
{\large\textbf{Simplify. Cancel the letters that are on the top and the bottom.}} {\color{muted}Show your working in each box.}\par\vspace{1mm}
\begin{tcbraster}[raster columns=2, raster column skip=4mm, raster row skip=3mm]
\qbox{1}{$6ab \div  2a$}{3.45cm}
\qbox{2}{$12xy \div  4y$}{3.45cm}
\qbox{3}{$10ab \div  5b$}{3.45cm}
\qbox{4}{$8xy \div  2x$}{3.45cm}
\qbox{5}{$14ab \div  7$}{3.45cm}
\qbox{6}{$x^{2} \div  x$}{3.45cm}
\qbox{7}{$a^{2} \div  a$}{3.45cm}
\qbox{8}{$4x^{2} \div  x$}{3.45cm}
\qbox{9}{$6x^{2} \div  2x$}{3.45cm}
\qbox{10}{$15a^{2} \div  5a$}{3.45cm}
\qbox{11}{$20x^{2} \div  4x$}{3.45cm}
\qbox{12}{$9x^{2} \div  3x$}{3.45cm}
\end{tcbraster}
\vfill\centerline{\color{good}\bfseries \ding{51}\ Check your answers at the back. Fix any you got wrong before moving on.}
\newpage
\banner{3}{Dividing Algebra}{Skill Check}{checkcol}
{\large Name: \fillline{8cm}}\hfill
\fcolorbox{checkcol}{white}{\parbox[c][9mm][c]{4.6cm}{\centering\color{checkcol}\bfseries\large Teacher mark: \underline{\hspace{9mm}} / 4}}
\par\vspace{2mm}
{\large\textbf{Simplify.}} Do this on your own --- no notes. Show your working, and write your final answer on the line.
\par\vspace{2mm}
\begin{tcbraster}[raster columns=2, raster column skip=5mm, raster row skip=5mm]
\cbox{1}{$12x \div  4$}{10.6cm}
\cbox{2}{$ab \div  b$}{10.6cm}
\cbox{3}{$10ab \div  5a$}{10.6cm}
\cbox{4}{$8x^{2} \div  2x$}{10.6cm}
\end{tcbraster}
\newpage
\banner{4}{Expanding Brackets}{Learn}{accent}
\begin{tcolorbox}[enhanced, colback=accentlight, colframe=accent, boxrule=0pt, leftrule=1.5mm, arc=0mm, top=1.5mm, bottom=1.5mm]\textbf{Goal:} I can expand brackets, e.g. $3(x + 2) = 3x + 6$ and $-2(x - 4) = -2x + 8$\end{tcolorbox}
\large

\textbf{Expanding} means getting rid of the brackets by multiplying.
A number or letter right in front of a bracket means \textbf{multiply everything inside}
the bracket by it.

\begin{stepsbox}{Multiply the outside by EVERY term inside}
Draw arrows so you don't miss one:
\[ 3(x + 2) \;\;\longrightarrow\;\; 3 \times x \;\text{ and }\; 3 \times 2 \;\;\longrightarrow\;\; 3x + 6 \]
\end{stepsbox}

\subsection*{Worked examples}
\begin{examples}
\example{$4(a + 5)$}{$4 \times  a = 4a$\\$4 \times  5 = 20$}{$= 4a + 20$}
\example{$2(x - 3)$}{$2 \times  x = 2x$\\$2 \times  (-3) = -6$}{$= 2x - 6$}
\example{$x(x + 5)$}{$x \times  x = x^{2}$\\$x \times  5 = 5x$}{$= x^{2} + 5x$}
\end{examples}

\subsection*{Negative numbers outside the bracket}
When the number outside is \textbf{negative}, use the sign rules:
\begin{center}
\renewcommand{\arraystretch}{1.3}
\begin{tabular}{|c|c|}
\hline
\rowcolor{accentlight} \textbf{Signs} & \textbf{Answer} \\ \hline
same signs: \; $+ \times +$ \; or \; $- \times -$ & \textbf{positive} $(+)$ \\ \hline
different signs: \; $+ \times -$ \; or \; $- \times +$ & \textbf{negative} $(-)$ \\ \hline
\end{tabular}
\end{center}

\begin{examples}
\example{$-2(x + 3)$}{$-2 \times  x = -2x$\\$-2 \times  3 = -6$}{$= -2x - 6$}
\example{$-3(x - 4)$}{$-3 \times  x = -3x$\\$-3 \times  (-4) = +12$}{$= -3x + 12$}
\example{$-(x - 5)$}{A minus on its own means $-1$\\$-1 \times  x = -x$,\; $-1 \times  (-5) = +5$}{$= -x + 5$}
\end{examples}

\begin{rememberbox}
The outside term multiplies \textbf{every} term inside: $3(x + 2)$ is $3x + 6$, \textbf{not} $3x + 2$.
With a negative outside, \textbf{every sign inside changes}: $-2(x - 4) = -2x + 8$.
\end{rememberbox}

\newpage\normalsize
\banner{4}{Expanding Brackets}{Practice A}{accent}
{\large\textbf{Expand.}} {\color{muted}Show your working in each box.}\par\vspace{1mm}
\begin{tcbraster}[raster columns=2, raster column skip=4mm, raster row skip=3mm]
\qbox{1}{$2(x + 3)$}{3.45cm}
\qbox{2}{$3(a + 4)$}{3.45cm}
\qbox{3}{$5(x + 1)$}{3.45cm}
\qbox{4}{$6(x + 2)$}{3.45cm}
\qbox{5}{$4(y - 2)$}{3.45cm}
\qbox{6}{$2(m - 5)$}{3.45cm}
\qbox{7}{$3(2x + 1)$}{3.45cm}
\qbox{8}{$2(4a + 3)$}{3.45cm}
\qbox{9}{$5(2y - 1)$}{3.45cm}
\qbox{10}{$x(x + 2)$}{3.45cm}
\qbox{11}{$a(a + 5)$}{3.45cm}
\qbox{12}{$x(x - 3)$}{3.45cm}
\end{tcbraster}
\vfill\centerline{\color{good}\bfseries \ding{51}\ Check your answers at the back. Fix any you got wrong before moving on.}
\newpage
\banner{4}{Expanding Brackets}{Practice B}{accent}
{\large\textbf{Expand. Be careful with the negative signs!}} {\color{muted}Show your working in each box.}\par\vspace{1mm}
\begin{tcbraster}[raster columns=2, raster column skip=4mm, raster row skip=3mm]
\qbox{1}{$-2(x + 3)$}{3.45cm}
\qbox{2}{$-3(a + 4)$}{3.45cm}
\qbox{3}{$-4(x - 2)$}{3.45cm}
\qbox{4}{$-5(y - 1)$}{3.45cm}
\qbox{5}{$-(x + 7)$}{3.45cm}
\qbox{6}{$-(a - 3)$}{3.45cm}
\qbox{7}{$2(3 - x)$}{3.45cm}
\qbox{8}{$4(2y - 3)$}{3.45cm}
\qbox{9}{$-2(3x - 1)$}{3.45cm}
\qbox{10}{$x(x - 6)$}{3.45cm}
\qbox{11}{$-x(x + 2)$}{3.45cm}
\qbox{12}{$-a(a - 4)$}{3.45cm}
\end{tcbraster}
\vfill\centerline{\color{good}\bfseries \ding{51}\ Check your answers at the back. Fix any you got wrong before moving on.}
\newpage
\banner{4}{Expanding Brackets}{Skill Check}{checkcol}
{\large Name: \fillline{8cm}}\hfill
\fcolorbox{checkcol}{white}{\parbox[c][9mm][c]{4.6cm}{\centering\color{checkcol}\bfseries\large Teacher mark: \underline{\hspace{9mm}} / 4}}
\par\vspace{2mm}
{\large\textbf{Expand.}} Do this on your own --- no notes. Show your working, and write your final answer on the line.
\par\vspace{2mm}
\begin{tcbraster}[raster columns=2, raster column skip=5mm, raster row skip=5mm]
\cbox{1}{$3(x + 5)$}{10.6cm}
\cbox{2}{$2(4a - 3)$}{10.6cm}
\cbox{3}{$x(x - 2)$}{10.6cm}
\cbox{4}{$-4(x - 3)$}{10.6cm}
\end{tcbraster}
\newpage
\banner{5}{Factorising into Brackets}{Learn}{accent}
\begin{tcolorbox}[enhanced, colback=accentlight, colframe=accent, boxrule=0pt, leftrule=1.5mm, arc=0mm, top=1.5mm, bottom=1.5mm]\textbf{Goal:} I can factorise, e.g. $2x + 6 = 2(x + 3)$\end{tcolorbox}
\large

\textbf{Factorising} is the \textbf{opposite of expanding}. We put the brackets back in.
\begin{center}
\colorbox{accentlight}{\begin{minipage}{0.8\linewidth}\centering\large\vspace{4pt}
$2(x + 3)$ \;$\xrightarrow{\text{ expand }}$\; $2x + 6$
\qquad\qquad
$2x + 6$ \;$\xrightarrow{\text{ factorise }}$\; $2(x + 3)$
\vspace{4pt}\end{minipage}}
\end{center}

\begin{stepsbox}{How to factorise}
\begin{enumerate}[nosep]
\item Find the \textbf{biggest number} (and/or \textbf{letter}) that goes into \textbf{every} term.
      Write it outside the bracket.
\item \textbf{Divide} each term by it. Write the answers inside the bracket.
\item \textbf{Check} by expanding --- you should get back to the question.
\end{enumerate}
\end{stepsbox}

\subsection*{Worked examples}
\begin{examples}
\example{$5a + 10$}{5 goes into both\\$5a \div  5 = a$,\; $10 \div  5 = 2$}{$= 5(a + 2)$}
\example{$6x - 9$}{3 goes into both\\$6x \div  3 = 2x$,\; $-9 \div  3 = -3$}{$= 3(2x - 3)$}
\example{$x^{2} + 4x$}{$x$ goes into both\\$x^{2} \div  x = x$,\; $4x \div  x = 4$}{$= x(x + 4)$}
\example{$3x^{2} + 6x$}{$3x$ goes into both\\$3x^{2} \div  3x = x$,\; $6x \div  3x = 2$}{$= 3x(x + 2)$}
\example{$14 - 7x$}{7 goes into both\\$14 \div  7 = 2$,\; $-7x \div  7 = -x$}{$= 7(2 - x)$}
\example{$8x - 12$}{4 is the \textbf{biggest}\\(2 works, but isn't the biggest)}{$= 4(2x - 3)$}
\end{examples}

\begin{rememberbox}
Keep the sign with each term: in $6x - 9$ the second term is $-9$, so the bracket has a minus: $3(2x - 3)$.
Always \textbf{check by expanding} your answer.
\end{rememberbox}

\newpage\normalsize
\banner{5}{Factorising into Brackets}{Practice A}{accent}
{\large\textbf{Factorise.}} {\color{muted}Show your working in each box.}\par\vspace{1mm}
\begin{tcbraster}[raster columns=2, raster column skip=4mm, raster row skip=3mm]
\qbox{1}{$3x + 6$}{3.45cm}
\qbox{2}{$2a + 8$}{3.45cm}
\qbox{3}{$5x + 15$}{3.45cm}
\qbox{4}{$7m + 7$}{3.45cm}
\qbox{5}{$4y - 12$}{3.45cm}
\qbox{6}{$2x - 10$}{3.45cm}
\qbox{7}{$6a + 9$}{3.45cm}
\qbox{8}{$10x + 15$}{3.45cm}
\qbox{9}{$x^{2} + 3x$}{3.45cm}
\qbox{10}{$a^{2} + 5a$}{3.45cm}
\qbox{11}{$x^{2} - 2x$}{3.45cm}
\qbox{12}{$y^{2} + y$}{3.45cm}
\end{tcbraster}
\vfill\centerline{\color{good}\bfseries \ding{51}\ Check your answers at the back. Fix any you got wrong before moving on.}
\newpage
\banner{5}{Factorising into Brackets}{Practice B}{accent}
{\large\textbf{Factorise fully. Take out the biggest factor you can.}} {\color{muted}Show your working in each box.}\par\vspace{1mm}
\begin{tcbraster}[raster columns=2, raster column skip=4mm, raster row skip=3mm]
\qbox{1}{$8x - 12$}{3.45cm}
\qbox{2}{$6a - 15$}{3.45cm}
\qbox{3}{$12y + 18$}{3.45cm}
\qbox{4}{$20x - 10$}{3.45cm}
\qbox{5}{$9m - 6$}{3.45cm}
\qbox{6}{$14 - 7x$}{3.45cm}
\qbox{7}{$x^{2} - 7x$}{3.45cm}
\qbox{8}{$ab + 3a$}{3.45cm}
\qbox{9}{$3x^{2} + 6x$}{3.45cm}
\qbox{10}{$2a^{2} - 8a$}{3.45cm}
\qbox{11}{$5y^{2} + 10y$}{3.45cm}
\qbox{12}{$4x^{2} - 6x$}{3.45cm}
\end{tcbraster}
\vfill\centerline{\color{good}\bfseries \ding{51}\ Check your answers at the back. Fix any you got wrong before moving on.}
\newpage
\banner{5}{Factorising into Brackets}{Skill Check}{checkcol}
{\large Name: \fillline{8cm}}\hfill
\fcolorbox{checkcol}{white}{\parbox[c][9mm][c]{4.6cm}{\centering\color{checkcol}\bfseries\large Teacher mark: \underline{\hspace{9mm}} / 4}}
\par\vspace{2mm}
{\large\textbf{Factorise fully.}} Do this on your own --- no notes. Show your working, and write your final answer on the line.
\par\vspace{2mm}
\begin{tcbraster}[raster columns=2, raster column skip=5mm, raster row skip=5mm]
\cbox{1}{$5x + 20$}{10.6cm}
\cbox{2}{$6a - 9$}{10.6cm}
\cbox{3}{$x^{2} + 8x$}{10.6cm}
\cbox{4}{$4y^{2} - 12y$}{10.6cm}
\end{tcbraster}
\newpage
\banner{$\star$}{Mixed Review}{Extra practice}{accent}
{\large\textbf{A mix of all five skills.}} Look carefully --- which skill do you need? {\color{muted}Answers are at the back.}\par\vspace{1mm}
\begin{tcbraster}[raster columns=3, raster column skip=4mm, raster row skip=3mm]
\qbox{1}{$5a + 2a$}{3.05cm}
\qbox{2}{$3 \times  2x$}{3.05cm}
\qbox{3}{$12y \div  3$}{3.05cm}
\qbox{4}{$2(x + 5)$}{3.05cm}
\qbox{5}{$3x + 12$}{3.05cm}
\qbox{6}{$4x + 3y + 2x$}{3.05cm}
\qbox{7}{$a \times  4b$}{3.05cm}
\qbox{8}{$10ab \div  2b$}{3.05cm}
\qbox{9}{$x(x + 1)$}{3.05cm}
\qbox{10}{$x^{2} + 7x$}{3.05cm}
\qbox{11}{$8m - 3m$}{3.05cm}
\qbox{12}{$5x \times  x$}{3.05cm}
\qbox{13}{$6x^{2} \div  x$}{3.05cm}
\qbox{14}{$4(2a - 1)$}{3.05cm}
\qbox{15}{$5y - 20$}{3.05cm}
\qbox{16}{$2x^{2} + 3x + x^{2}$}{3.05cm}
\qbox{17}{$4a \times  2a$}{3.05cm}
\qbox{18}{$16xy \div  8x$}{3.05cm}
\qbox{19}{$-3(x - 2)$}{3.05cm}
\qbox{20}{$6x^{2} + 9x$}{3.05cm}
\end{tcbraster}
\newpage
{\LARGE\bfseries\color{accent} Answers to the Practice Questions}\par\vspace{1mm}
{\small\color{muted} Terms can be in a different order and still be right --- e.g. $2b + 3a$ is the same as $3a + 2b$.}\par
\subsection*{1. Collecting Like Terms}
\textbf{\color{muted}Practice A}
\begin{multicols}{4}\begin{enumerate}[nosep, label=\arabic*., start=1]
\item $2a$
\item $3a$
\item $5a$
\item $9x$
\item $4m$
\item $4b$
\item $5y$
\item $6x$
\item $6p$
\item $6c$
\item $2x^{2}$
\item $7x^{2}$
\end{enumerate}\end{multicols}

\textbf{\color{muted}Practice B}
\begin{multicols}{4}\begin{enumerate}[nosep, label=\arabic*., start=1]
\item $6a + 3b$
\item $6x + 2y$
\item $4a + 6b$
\item $6x + 3$
\item $5m + 3n$
\item $2a + 2b$
\item $4x + 5$
\item $5x + 2y$
\item $6ab$
\item $x^{2} + 5x$
\item $3x^{2} + x$
\item $4y^{2} + 6y$
\end{enumerate}\end{multicols}

\subsection*{2. Multiplying Algebra}
\textbf{\color{muted}Practice A}
\begin{multicols}{4}\begin{enumerate}[nosep, label=\arabic*., start=1]
\item $4a$
\item $5x$
\item $ab$
\item $xy$
\item $6a$
\item $12x$
\item $15m$
\item $12x$
\item $a^{2}$
\item $y^{2}$
\item $2ab$
\item $pqr$
\end{enumerate}\end{multicols}

\textbf{\color{muted}Practice B}
\begin{multicols}{4}\begin{enumerate}[nosep, label=\arabic*., start=1]
\item $6ab$
\item $8xy$
\item $10ab$
\item $6ab$
\item $3a^{2}$
\item $10x^{2}$
\item $12m^{2}$
\item $6x^{2}$
\item $6ab$
\item $x^{3}$
\item $x^{3}$
\item $6x^{3}$
\end{enumerate}\end{multicols}

\subsection*{3. Dividing Algebra}
\textbf{\color{muted}Practice A}
\begin{multicols}{4}\begin{enumerate}[nosep, label=\arabic*., start=1]
\item $4a$
\item $4x$
\item $2b$
\item $5y$
\item $3x$
\item $3p$
\item $2x$
\item $m$
\item $a$
\item $y$
\item $5$
\item $7$
\end{enumerate}\end{multicols}

\textbf{\color{muted}Practice B}
\begin{multicols}{4}\begin{enumerate}[nosep, label=\arabic*., start=1]
\item $3b$
\item $3x$
\item $2a$
\item $4y$
\item $2ab$
\item $x$
\item $a$
\item $4x$
\item $3x$
\item $3a$
\item $5x$
\item $3x$
\end{enumerate}\end{multicols}

\subsection*{4. Expanding Brackets}
\textbf{\color{muted}Practice A}
\begin{multicols}{4}\begin{enumerate}[nosep, label=\arabic*., start=1]
\item $2x + 6$
\item $3a + 12$
\item $5x + 5$
\item $6x + 12$
\item $4y - 8$
\item $2m - 10$
\item $6x + 3$
\item $8a + 6$
\item $10y - 5$
\item $x^{2} + 2x$
\item $a^{2} + 5a$
\item $x^{2} - 3x$
\end{enumerate}\end{multicols}

\textbf{\color{muted}Practice B}
\begin{multicols}{4}\begin{enumerate}[nosep, label=\arabic*., start=1]
\item $-2x - 6$
\item $-3a - 12$
\item $-4x + 8$
\item $-5y + 5$
\item $-x - 7$
\item $-a + 3$
\item $6 - 2x$
\item $8y - 12$
\item $-6x + 2$
\item $x^{2} - 6x$
\item $-x^{2} - 2x$
\item $-a^{2} + 4a$
\end{enumerate}\end{multicols}

\subsection*{5. Factorising into Brackets}
\textbf{\color{muted}Practice A}
\begin{multicols}{4}\begin{enumerate}[nosep, label=\arabic*., start=1]
\item $3(x + 2)$
\item $2(a + 4)$
\item $5(x + 3)$
\item $7(m + 1)$
\item $4(y - 3)$
\item $2(x - 5)$
\item $3(2a + 3)$
\item $5(2x + 3)$
\item $x(x + 3)$
\item $a(a + 5)$
\item $x(x - 2)$
\item $y(y + 1)$
\end{enumerate}\end{multicols}

\textbf{\color{muted}Practice B}
\begin{multicols}{4}\begin{enumerate}[nosep, label=\arabic*., start=1]
\item $4(2x - 3)$
\item $3(2a - 5)$
\item $6(2y + 3)$
\item $10(2x - 1)$
\item $3(3m - 2)$
\item $7(2 - x)$
\item $x(x - 7)$
\item $a(b + 3)$
\item $3x(x + 2)$
\item $2a(a - 4)$
\item $5y(y + 2)$
\item $2x(2x - 3)$
\end{enumerate}\end{multicols}

\subsection*{$\star$ Mixed Review}
\textbf{\color{muted}}
\begin{multicols}{4}\begin{enumerate}[nosep, label=\arabic*., start=1]
\item $7a$
\item $6x$
\item $4y$
\item $2x + 10$
\item $3(x + 4)$
\item $6x + 3y$
\item $4ab$
\item $5a$
\item $x^{2} + x$
\item $x(x + 7)$
\item $5m$
\item $5x^{2}$
\item $6x$
\item $8a - 4$
\item $5(y - 4)$
\item $3x^{2} + 3x$
\item $8a^{2}$
\item $2y$
\item $-3x + 6$
\item $3x(2x + 3)$
\end{enumerate}\end{multicols}

\end{document}
